%% ma: L12-B08-0187
%% lop: 12
%% chuong: 2
%% bai: 8
%% dang: 12.8.4
%% loai: TN4
%% muc: VD
%% dap_an: "A"
%% nguon: "CĐ2.6.pdf giáo viên gửi 27/10/2026 (M07, Lê Minh Khôi – Nguyễn Anh Khôi, Chương II Chủ đề 6 Ứng dụng vào các bài toán thực tế), Câu 15"
%% kien_thuc: "Bài 8 (tọa độ vectơ, khoảng cách), chuyển động thẳng đều, lớp 10 (bất phương trình bậc hai)"
%% trang_thai: chinh-thuc
\begin{ex}
Trong không gian $Oxyz$ mỗi đơn vị trên hệ trục ứng với $10$ km, trạm kiểm soát không lưu đang theo dõi hai máy bay. Máy bay thứ nhất ban đầu ở tọa độ $A(25;-10;1)$ và bay theo hướng vectơ $\vec v_1=(-3;-4;0)$ với tốc độ không đổi là $750$ km/h. Máy bay thứ hai ban đầu ở tọa độ $B(30;-25;1{,}1)$ và bay theo hướng vectơ $\vec v_2=(-4;3;0)$ với tốc độ không đổi là $900$ km/h. Trên máy bay thứ nhất có gắn ra đa tránh va chạm với bán kính hoạt động $50$ km. Hỏi thời gian máy bay thứ hai xuất hiện trên màn hình ra đa của máy bay thứ nhất là bao nhiêu phút (làm tròn kết quả đến hàng phần trăm)?
\choice{\True $4{,}91$ phút}{$4{,}79$ phút}{$8{,}23$ phút}{$5{,}51$ phút}
\loigiai{Sau $t$ giờ (đơn vị trục): $M_1=A+75t\cdot\dfrac{\vec v_1}5=(25-45t;-10-60t;1)$, $M_2=B+90t\cdot\dfrac{\vec v_2}5=(30-72t;-25+54t;1{,}1)$.

$M_1M_2^2=(5-27t)^2+(-15+114t)^2+0{,}01\le25\Leftrightarrow t_1\le t\le t_2$ với $t_1\approx0{,}0935$, $t_2\approx0{,}1754$.

Thời gian $(t_2-t_1)\cdot60\approx4{,}91$ phút.}
\end{ex}
